% ============================================================================
% ARBEITSBLATT DS 6 - Gruppe A: Personenbeschreibung (Brückenthema)
% ============================================================================

\documentclass[11pt,a4paper]{article}

\usepackage[utf8]{inputenc}
\usepackage[T1]{fontenc}
\usepackage[ngerman]{babel}
\usepackage[left=2cm,right=2cm,top=1.8cm,bottom=1.5cm]{geometry}
\usepackage{helvet}
\usepackage[table]{xcolor}
\usepackage{tabularx}
\usepackage{array}
\usepackage{enumitem}
\usepackage{tcolorbox}
\usepackage{fancyhdr}
\usepackage{lastpage}
\usepackage{amssymb}

\renewcommand{\familydefault}{\sfdefault}

\definecolor{spanisch}{HTML}{c41e3a}
\definecolor{grau}{HTML}{666666}
\definecolor{hellgrau}{HTML}{f5f5f5}

\pagestyle{fancy}
\fancyhf{}
\fancyhead[L]{\small\textcolor{grau}{Spanisch Spätbeginner -- Gruppe A}}
\fancyhead[R]{\small\textcolor{grau}{AB-06-A}}
\fancyfoot[C]{\small\thepage/\pageref{LastPage}}
\renewcommand{\headrulewidth}{0.5pt}

\newcommand{\luecke}[1]{\underline{\hspace{#1}}}
\newcommand{\punktebox}[1]{\hfill\fbox{\small \_\_/#1}}
\newcommand{\es}[1]{\textit{#1}}

\newtcolorbox{merkkasten}[1][]{
    colback=white,
    colframe=spanisch,
    fonttitle=\bfseries\color{spanisch},
    title=#1,
    boxrule=1.5pt,
    arc=0pt,
    left=5pt, right=5pt, top=5pt, bottom=5pt
}

\newtcolorbox{buchverweis}{
    colback=hellgrau,
    colframe=grau,
    boxrule=0.5pt,
    arc=0pt,
    left=5pt, right=5pt, top=3pt, bottom=3pt
}

\newcounter{aufgabe}
\newcommand{\aufgabe}[2]{%
    \stepcounter{aufgabe}%
    \vspace{0.8em}%
    \noindent\textbf{\theaufgabe.} #1 \punktebox{#2}%
    \vspace{0.3em}%
}

\begin{document}

% ============================================================================
% KOPFBEREICH
% ============================================================================
\noindent
\begin{tabularx}{\textwidth}{@{}Xr@{}}
    {\Large\textbf{\textcolor{spanisch}{Personas y características}}} & Name: \luecke{5cm} \\[0.3em]
    {\small DS 6 -- Brückenthema} & Datum: \luecke{3cm} \\
\end{tabularx}

\vspace{0.3em}
\hrule
\vspace{0.8em}

% ============================================================================
% MERKKASTEN 1: ser vs. estar + Adjektiv
% ============================================================================
\begin{merkkasten}[Kasten 1: ser vs. estar + Adjektiv]
\vspace{0.2em}

\begin{tabularx}{\textwidth}{@{}|p{2.5cm}|X|X|@{}}
\hline
\rowcolor{hellgrau}
\textbf{Verb} & \textbf{Verwendung} & \textbf{Beispiel} \\
\hline
\textbf{ser} & \luecke{3.5cm} (dauerhaft) & Mi hermana \luecke{1.5cm} simpática. \\[0.6em]
\hline
\textbf{estar} & \luecke{3.5cm} (momentan) & Hoy \luecke{1.5cm} cansado. \\[0.6em]
\hline
\end{tabularx}

\vspace{0.3em}
\textbf{Merkhilfe:} \es{Ser} = so bin ich IMMER $\bullet$ \es{Estar} = so bin ich JETZT
\vspace{0.2em}
\end{merkkasten}

\vspace{0.8em}

% ============================================================================
% MERKKASTEN 2: Konnektoren
% ============================================================================
\begin{merkkasten}[Kasten 2: Konnektoren für Personenbeschreibungen]
\vspace{0.2em}

\begin{tabularx}{\textwidth}{@{}|p{3cm}|X|@{}}
\hline
\rowcolor{hellgrau}
\textbf{Konnektor} & \textbf{Beispiel} \\
\hline
\textbf{también} (auch) & Mi padre es trabajador. También es muy \luecke{2.5cm}. \\[0.4em]
\hline
\textbf{además} (außerdem) & Es simpática. Además, es muy \luecke{2.5cm}. \\[0.4em]
\hline
\textbf{pero} (aber) & Es inteligente, pero a veces es \luecke{2.5cm}. \\[0.4em]
\hline
\textbf{porque} (weil) & Me gusta mucho porque es muy \luecke{2.5cm}. \\[0.4em]
\hline
\end{tabularx}

\vspace{0.3em}
\end{merkkasten}

\vspace{1em}

% ============================================================================
% AUFGABEN
% ============================================================================

\aufgabe{Adjektive: Ordne zu und ergänze die Endung (-o/-a).}{8}

\begin{tabularx}{\textwidth}{@{}|X|X|X|X|@{}}
\hline
\rowcolor{hellgrau}
\textbf{Spanisch} & \textbf{Deutsch} & \textbf{Spanisch} & \textbf{Deutsch} \\
\hline
simpátic\_\_ & \luecke{2.5cm} & paciente & \luecke{2.5cm} \\[0.5em]
\hline
inteligente & \luecke{2.5cm} & amable & \luecke{2.5cm} \\[0.5em]
\hline
trabajador\_\_ & \luecke{2.5cm} & divertid\_\_ & \luecke{2.5cm} \\[0.5em]
\hline
responsable & \luecke{2.5cm} & creativ\_\_ & \luecke{2.5cm} \\[0.5em]
\hline
\end{tabularx}

\vspace{0.8em}

\aufgabe{Ergänze \es{ser} oder \es{estar} in der richtigen Form.}{6}

\begin{enumerate}[label=\alph*), leftmargin=2em]
    \item Mi madre \luecke{2cm} muy paciente. (Eigenschaft)
    \item Hoy mis hermanos \luecke{2cm} cansados. (Zustand)
    \item ¿Cómo \luecke{2cm} tu mejor amigo? -- Es muy divertido.
    \item Nosotros \luecke{2cm} trabajadores.
    \item María \luecke{2cm} enferma hoy. (momentan)
    \item Yo \luecke{2cm} responsable y organizada.
\end{enumerate}

\vspace{0.8em}

\aufgabe{Beschreibe eine Person deiner Wahl (Familie/Freund). Nutze mindestens 3 Adjektive und 2 Konnektoren.}{8}

\begin{buchverweis}
\textbf{Beispiel:} \es{Mi mejor amiga se llama Laura. Es muy simpática y también es inteligente. Además, es muy divertida, pero a veces es un poco impaciente.}
\end{buchverweis}

\vspace{0.5em}
\begin{tabularx}{\textwidth}{@{}|X|@{}}
\hline
\\[0.5em]
\luecke{14cm} \\[0.8em]
\luecke{14cm} \\[0.8em]
\luecke{14cm} \\[0.8em]
\luecke{14cm} \\[0.5em]
\hline
\end{tabularx}

\vspace{1em}

\aufgabe{Partnerdialog: Interviewe deinen Partner über seine Familie.}{6}

\begin{buchverweis}
\textbf{Fragen:} \es{¿Cómo eres? / ¿Cómo es tu mejor amigo/-a? / ¿Cómo es tu familia? / ¿Por qué?}
\end{buchverweis}

\vspace{0.3em}
\textbf{Notiere 3 Informationen über deinen Partner:}

\vspace{0.3em}
1. \luecke{12cm}

\vspace{0.3cm}
2. \luecke{12cm}

\vspace{0.3cm}
3. \luecke{12cm}

\vspace{1em}

\aufgabe{Schreibe eine Zusammenfassung über deinen Partner (Overlap-Vorbereitung).}{5}

\begin{buchverweis}
\textbf{Struktur:} \es{[Nombre] es ... y ... . También es ... . Su mejor amigo/-a es ... .}
\end{buchverweis}

\vspace{0.5em}
\begin{tabularx}{\textwidth}{@{}|X|@{}}
\hline
\\[0.5em]
\luecke{14cm} \\[0.8em]
\luecke{14cm} \\[0.8em]
\luecke{14cm} \\[0.5em]
\hline
\end{tabularx}

\vspace{1em}
\hrule
\vspace{0.3em}
\begin{flushright}
\textbf{Gesamt: \luecke{1cm} / 33 Punkte}
\end{flushright}

\end{document}
